\section{Introduction}\label{sec:intro}

Ask a language model to place 21 circles in a unit square to maximize the sum of their radii.
It will not search. It reaches for the nearest square grid, five by five ($k = 5$) at
radius $1/(2k)$, and truncates it, leaving four cells empty. A companion submission
\citep{companion} gives that behaviour a closed form, naming the modal output in advance.
It reports that form as the empirical mode at all seven $N$ tested, at four of the five
held out, and under a perturbed container and paraphrase. The two cell sets are listed in
Table~\ref{tab:glossary}. This paper's own generation-0 rows were sampled
fresh under its loop arm, on the same instrument and prompt bytes
\citep{companion}. At $N = 31$ they put the anchor on 7 of 8 valid
outputs. At $N = 13$, on 4 of 9 (Table~\ref{tab:overview}).

A characterization that sharp invites an inference. What a proposer emits zero-shot is what it
will contribute inside a discovery loop, so the loop need not run. Discovery systems in the
FunSearch line \citep{romeraparedes2024funsearch,novikov2025alphaevolve} do not call their
proposer zero-shot. They condition it on a parent and iterate it against an archive. Each is a
change of setting, as is the serving path \citep{companion}. Each can break the
characterization.

\textbf{Objects and arms.} A \emph{template} is one
member of the grid-plus-filler recipe family. The \emph{anchor} is the template the model's
outputs concentrate on at a given $N$. The \emph{family argmax} is the largest value the
family can express at that $N$. A \emph{discriminating cell} is an $N$ where the two differ,
so only there can anchoring be told from family search. Table~\ref{tab:glossary} carries the
rest. Eight registered arms change four things. Arm MU changes
the parent. Arm L adds iteration, feeding a scored proposal back as the next parent. Arms P
and P-D change the serving path and the reasoning setting on it. Arms GM, GM2, GM3 and V
change the vendor, GM2 and GM3 differing only by output budget.

\textbf{Claim.} A weak-tier LLM proposer's zero-shot anchor in circle packing is contingent on
the serving path and its request parameters. A zero-shot
characterization therefore does not establish that its anchor transfers into a discovery
loop, absent a measurement there. The anchor does not transfer. The family argmax does, as a
ceiling. Under the clearance rule of \S\ref{sec:method}, no output valid at $10^{-9}$
exceeds it by more than $10^{-6}$ in any arm. At the second vendor's discriminating cells
most valid outputs sit on it. Two kinds of residual are reported and neither is a clearance.
The primary $10^{-6}$ tolerance admits circle overlaps up to
that size, and under it 18 arm MU outputs sit above the argmax by at most
$1.53\times10^{-6}$; all 18 fail at $10^{-9}$
(\S\ref{sec:arm-mu}). In arm GM3, 29 outputs valid at $10^{-9}$ sit above the argmax by at
most $5.3\times10^{-10}$, the rounding of a printed coordinate (\S\ref{sec:gm3-v}). Two arms carry no row in Table~\ref{tab:overview}. Arm
GM finished UNDERPOWERED whether its routed rows are counted or not. Arm GM2 returned nothing
parseable at the smaller output budget. Arm GM3 diagnosed the cause at the larger.

\begin{table}[!ht]
\caption{Six of the eight arms: what each changed, its bar, the outcome
and the verdict. Codes like P-MU1 name a registered prediction,
written out in its arm's section and Supplement~S1. L1 to L4 are the loop arm's
four, named in \S\ref{sec:arm-l}. Arms MU and L ran on the agent-runtime instrument, P and
P-D on the pinned path. Bracketed ranges are 95\% Wilson intervals on the rate they follow, reported as descriptive precision; no verdict turns on one.}
\label{tab:overview}
\begin{center}
\footnotesize
\setlength{\tabcolsep}{3pt}
\begin{tabular}{@{}>{\raggedright\arraybackslash}p{0.05\linewidth}>{\raggedright\arraybackslash}p{0.15\linewidth}>{\raggedright\arraybackslash}p{0.24\linewidth}>{\raggedright\arraybackslash}p{0.22\linewidth}>{\raggedright\arraybackslash}p{0.21\linewidth}@{}}
\toprule
Arm & What changed & Registered bar & Outcome & Verdict \\
\midrule
MU & The parent the proposer is conditioned on & Under a better in-family parent, on-prediction rate
at most 20\% and keep-or-improve at least 50\% & 26 of 26 keep or improve [87--100\%], 25 within $10^{-6}$ of the parent's score and one
$1.2\times10^{-6}$ above it, so the on-prediction rate there is 0 of 26 [0--13\%]; under the model's
own anchor, 7 of 28 keep it [13--43\%] and 17 move above it & The rule reduces to one
condition under the better parent; the own-anchor rate misses P-MU1's 50\% bar, so the
anchor is not self-stable; the falsifier is that rule's negation and adds no second result \\
\addlinespace
L & Five archive-and-mutate generations & Pooled conditioned on-prediction rate below the
generation-0 rate in both cells (L2); the falsifier is its negation & $N = 13$: 4 of 9 at generation 0, 5 of
20 conditioned. $N = 31$: 7 of 8, then 6 of 29. Pooled 11 of 17 [41--83\%], then 11 of 49
[13--36\%] & L2 holds; L1 and L3 not met; L4 open (0 of 49 clear the family, [0--7\%]) \\
\addlinespace
P & The serving path: the bare pinned chat-completions endpoint & A floor of five valid
outputs per cell & 0 of 105 valid at the square cells [0--4\%]; 4 of 75 at the held-out cells
[2--13\%] & Every direct-emission prediction UNSCOREABLE: no cell reaches the floor \\
\addlinespace
P-D & Extended thinking on against off, same pinned path & The reasoning budget is the cause,
against the missing context, against anything else & Thinking on, 12 of 14
valid with the anchor the plurality value at 5 of 12 [19--68\%]; off, 0 of 13 [0--23\%] &
The reasoning-budget hypothesis holds; extended thinking is the component \\
\addlinespace
GM3 & A second vendor's gemma at a 16,384-token budget & Pooled rate at least 30\%, and the
anchor modal in 5 of 5 scoreable cells & 46 of 80 pooled (57.5\%, [47--68\%]); modal
in 2 of 5; 11 of 43 at the discriminating cells (26\%, [15--40\%]) & The pooled bar is met;
the discriminating cells, the only ones that test anchoring, miss it (post hoc split) \\
\addlinespace
V & Thirteen free-tier aliases & At least 10 of a model's 25 invocations valid, then V1 at
least 50\% and V2 at most 10\% & Ten aliases below the floor; two aliases meet the V1 pooled
bar at point estimate, on 8 of 12 [39--86\%] and 11 of 18 [39--80\%] of their valid outputs,
and a third does not, at 0 of 15 [0--20\%] & The boundary is located, not attributed \\
\bottomrule
\end{tabular}
\end{center}
\end{table}

\textbf{Contributions.}
\begin{enumerate}
\item \emph{A measurement of transfer, not an argument about it.} Each change of setting got
its own arm, with thresholds and a falsifier fixed before the first
invocation. Every deviation is tabled.
\item \emph{A named instrument component, at one cell.} On the pinned path at $N = 13$, arm
P's validity collapse survives the one system line tested, a request for care. One
registered condition (P-D, whose rows both submissions report) and two post hoc diagnostics
separate the candidates. Extended thinking carries the anchor.
\item \emph{A cross-vendor boundary that is located but not attributed.} Two second-vendor
measurements disagree, with model variant and serving path confounded.
\end{enumerate}

\textbf{Scope.} The proposer is one vendor's weak-tier alias
(\texttt{claude-haiku-4.5}) except where another is named. The conditioning arm runs
at three trap cells, $N = 13$, 21 and 31. The loop arm runs at two of them.

\emph{What is not claimed.} Nothing about the companion paper's numerical baseline, about
which constructions beat the family argmax, or about its
mid-tier alias. Those are its results, cited and not restated \citep{companion}.
Nothing about mechanism inside the weights, about what a discovery system does at
generation zero, or whether diversity machinery
\citep{mouret2015mapelites,lehman2011novelty} repairs or explores. Preregistration
is an adopted standard, not a contribution (\S\ref{sec:related}).
